\begin{encadre}
\textbf{Réponse courte.}
\begin{enumerate}[nosep]
\item \textbf{Oui près de la surface.} Sous l'empreinte hétérogène du TFE, la déformation
      principale majeure maximale est multipliée par 1,6 à 2,0 à \SI{1}{mm}
      de profondeur, par 1,5 à 1,8 à \SI{2}{mm} et par 1,3 à 1,5 à \SI{5}{mm},
      selon la texture (rainurage FAA, macrotextures de MPD 0,6 à \SI{1.5}{mm}). Les
      cisaillements proches de la surface sont amplifiés davantage encore.
\item \textbf{Non en profondeur.} Le facteur tombe à 1,05--1,23 à \SI{10}{mm} et à
      1,00--1,01 à \SI{20}{mm} ; l'effet est nul à \SI{40}{mm} (1,00--1,00). La déformation en base de couche liée et le cisaillement
      maximal vers \SI{100}{mm}, qui sont les deux grandeurs clés du TFE, ne sont pas modifiés.
\item La texture crée donc une \textbf{seconde zone critique, superficielle} (0 à
      \SI{10}{mm}), distincte de celle du TFE. C'est la zone d'amorçage de la fissuration par le
      haut et de l'arrachement.
\item L'effet croît avec la MPD, mais \textbf{la MPD ne suffit pas} : à MPD égale, la pente des
      aspérités, l'asymétrie de la texture et surtout la raideur du caoutchouc changent le
      résultat.
\item \textbf{Confiance} : les tendances et les profondeurs d'influence sont solides ; la valeur
      absolue aux premiers millimètres est incertaine d'un facteur de l'ordre de 1,5, à cause de
      paramètres mal connus (\S\ref{sec:verif}).
\end{enumerate}
\end{encadre}